%% kapa: exam=1 round=32 year=2021 period=1 form=A subject=real_estate subject_ko=부동산학원론 seq=14 number=94
%% source: Q-Net 감정평가사 기출문제 (cst00309) artlSeq=5208397 file=extracted/1008/5208397/A형/2021년도 제32회 감정평가사 1차 1교시 A형.pdf
%% conversion: kordoc 4.24.51 -> tools/gen1.py (auto); review=proofread-sonnet; flags=none
\begin{question}{94}
다음 중 부동산시장과 부동산정책에 관한 설명으로 옳은 것은 몇 개인가?
\begin{framedtext}
○ 부동산정책이 자원배분의 비효율성을 오히려 악화시키는 것을 시장의 실패라 한다. \par
○ 법령상 도입순서를 비교하면 부동산거래신고제는 부동산실명제보다 빠르다. \par
○ 개발행위허가제와 택지소유상한제는 현재 시행되고 있는 제도이다. \par
○ 분양가상한제와 개발부담금제는 정부가 직접적으로 부동산시장에 개입하는 정책수단이다. \par
○ PIR(Price to Income Ratio)은 가구의 주택지불능력을 측정하는 지표이다. \par
\end{framedtext}
\begin{choices}
  \item 없음
  \item 1개
  \item 2개
  \item 3개
  \item 4개
\end{choices}
\end{question}
\answer{2}
